\section{AIME I 1986}
        \begin{enumerate}
        	\item (\textit{AIME I 1986}) What is the sum of the solutions to the equation $\sqrt[4]{x} = \frac{12}{7 - \sqrt[4]{x}}$ ?


 
 \textbf{Answer: }337



\textbf{Solution}

Let $y = \sqrt[4]{x}$. Then we have\textbf{$y(7 - y) = 12$}, or, by simplifying,\textbf{
 \[y^2 - 7y + 12 = (y - 3)(y - 4) = 0.\]}

 This means that $\sqrt[4]{x} = y = 3$ or \textbf{$4$}.

 Thus the sum of the possible solutions for \textbf{$x$} is \textbf{$4^4 + 3^4 = \boxed{337}$}.

 




	\item (\textit{AIME I 1986}) Evaluate the product 
 \[\left(\sqrt{5}+\sqrt{6}+\sqrt{7}\right)\left(\sqrt{5}+\sqrt{6}-\sqrt{7}\right)\left(\sqrt{5}-\sqrt{6}+\sqrt{7}\right)\left(-\sqrt{5}+\sqrt{6}+\sqrt{7}\right).\]


 
 \textbf{Answer: }104



\textbf{Solution 1 (Algebra: Generalized)}

More generally, let $(x,y,z)=\left(\sqrt5,\sqrt6,\sqrt7\right)$ so that $\left(x^2,y^2,z^2\right)=(5,6,7).$

 We rewrite the original expression in terms of $x,y,$ and $z,$ then apply the difference of squares repeatedly:
 \begin{align*} (x+y+z)(x+y-z)(x-y+z)(-x+y+z) &= \left[((x+y)+z)((x+y)-z)\right]\left[((z+(x-y))(z-(x-y))\right] \\ &= \left[(x+y)^2-z^2\right]\left[z^2 - (x-y)^2\right] \\ &= \left[x^2+2xy+y^2-z^2\right]\left[z^2-x^2+2xy-y^2\right] \\ &= \left[2xy + \left(x^2+y^2-z^2\right)\right]\left[2xy - \left(x^2+y^2-z^2\right)\right] \\ &= \left(2xy\right)^2 - \left(x^2+y^2-z^2\right)^2 \\ &= \left(2\cdot\sqrt5\cdot\sqrt6\right)^2 - \left(5+6-7\right)^2 \\ &= \boxed{104}. \end{align*}
<u>\textbf{Remark}</u>

 From this solution, note that the original expression has cyclic symmetry with respect to $x,y,$ and $z:$
 \begin{align*} (x+y+z)(x+y-z)(x-y+z)(-x+y+z) &= \cdots \\ &= \left(2xy\right)^2 - \left(x^2+y^2-z^2\right)^2 \\ &= 4x^2y^2 - x^4 - y^4 - z^4 - 2x^2y^2 + 2y^2z^2 + 2z^2x^2 \\ &= 2x^2y^2 + 2y^2z^2 + 2z^2x^2 - x^4 - y^4 - z^4. \end{align*}
~MRENTHUSIASM

 





\textbf{Solution 2 (Algebra: Specific)}

We group the first and last factors as well as the two middle factors, then apply the difference of squares repeatedly:
 \begin{align*} \left(\left(\sqrt{6} + \sqrt{7}\right)^2 - \sqrt{5}^2\right)\left(\sqrt{5}^2 - \left(\sqrt{6} - \sqrt{7}\right)^2\right) &= \left(13 + 2\sqrt{42} - 5\right)\left(5 - \left(13 - 2\sqrt{42}\right)\right) \\ &= \left(2\sqrt{42} + 8\right)\left(2\sqrt{42} - 8\right) \\ &= \left(2\sqrt{42}\right)^2 - 8^2 \\ &= \boxed{104}. \end{align*}
~Azjps (Solution)

 





\textbf{Solution 3 (Geometry)}

Notice that in a triangle with side-lengths $2\sqrt5,2\sqrt6,$ and $2\sqrt7,$ by Heron's Formula, the area is the square root of the original expression.

 Let $\theta$ be the measure of the angle opposite the $2\sqrt7$ side. By the Law of Cosines,
 \[\cos\theta=\frac{\left(2\sqrt5\right)^2+\left(2\sqrt{6}\right)^2-\left(2\sqrt7\right)^2}{2\cdot 2\sqrt5\cdot2\sqrt6}=\frac{16}{8\sqrt{30}}=\sqrt{\frac{2}{15}},\]so $\sin\theta=\sqrt{1-\cos^2\theta}=\sqrt{\frac{13}{15}}.$

 The area of the triangle is then
 \[\frac{\sin\theta}{2}\cdot 2\sqrt5\cdot 2\sqrt6=\sqrt{\frac{26}{30}}\cdot 2\sqrt{30}=\sqrt{104},\]so our answer is $\left(\sqrt{104}\right)^2=\boxed{104}.$

 




	\item (\textit{AIME I 1986}) If $\tan x+\tan y=25$ and $\cot x + \cot y=30$, what is $\tan(x+y)$?


 
 \textbf{Answer: }150



\textbf{Solution 1}

Since $\cot$ is the reciprocal function of $\tan$:

 $\cot x + \cot y = \frac{1}{\tan x} + \frac{1}{\tan y} = \frac{\tan x + \tan y}{\tan x \cdot \tan y} = 30$

 Thus, $\tan x \cdot \tan y = \frac{\tan x + \tan y}{30} = \frac{25}{30} = \frac{5}{6}$

 Using the tangent addition formula: 

 $\tan(x+y) = \frac{\tan x + \tan y}{1-\tan x \cdot \tan y} = \frac{25}{1-\frac{5}{6}} = \boxed{150}$.

 





\textbf{Solution 2}

Using the formula for tangent of a sum, $\tan(x+y)=\frac{\tan x + \tan y}{1-\tan x \tan y} = \frac{25}{1-\tan x \tan y}$. We only need to find $\tan x \tan y$.

 We know that $25 = \tan x + \tan y = \frac{\sin x}{\cos x} + \frac{\sin y}{\cos y}$. Cross multiplying, we have $\frac{\sin x \cos y + \cos x \sin y}{\cos x \cos y} = \frac{\sin(x+y)}{\cos x \cos y} = 25$.

 Similarly, we have $30 = \cot x + \cot y = \frac{\cos x}{\sin x} + \frac{\cos y}{\sin y} = \frac{\cos x \sin y + \sin x \cos y}{\sin x \sin y} = \frac{\sin(x+y)}{\sin x \sin y}$.

 Dividing:

 $\frac{25}{30} = \frac{\frac{\sin(x+y)}{\cos x \cos y}}{\frac{\sin(x+y)}{\sin x \sin y}} = \frac{\sin x \sin y}{\cos x \cos y} = \tan x \tan y = \frac{5}{6}$. Plugging in to the earlier formula, we have $\tan(x+y) = \frac{25}{1-\frac{5}{6}} = \frac{25}{\frac{1}{6}} = \boxed{150}$.

 





\textbf{Solution 3 (less trig required, use of quadratic formula)}

Let $a=\tan x$ and $b=\tan y$. This simplifies the equations to:

 
 \[a + b = 25\]

 
 \[\frac{1}{a} + \frac{1}{b} = 30\]

 Taking the tangent of a sum formula from Solution 2, we get $\tan(x+y) = \frac{25}{1 - ab}$. 

 We can use substitution to solve the system of equations from above: $b = -a + 25$, so $\frac{1}{a} + \frac{1}{-a + 25} = 30$. 

 Multiplying by $-a(a-25)$,  we get $a + (-a + 25) = -30a(a-25)$, which is $-30a^2 + 750a = 25$. Dividing everything by 5 and shifting everything to one side gives $6a^2 - 150a + 5 = 0$.

 Using the quadratic formula gives $a = \frac{150 \pm \sqrt {22380}}{12}$. Since this looks too hard to simplify, we can solve for $b$ using $a + b = 25$, which turns out to also be $b = \frac{150 \pm \sqrt {22380}}{12}$, provided that the sign of the radical in $a$ is opposite the one in $b$.

 WLOG, assume $a = \frac{150 + \sqrt{22380}}{12}$ and $b = \frac{150 - \sqrt{22380}}{12}$. Multiplying them gives $ab = \frac{22500 - 22380}{144}$ which simplifies to $\frac{5}{6}$.

 THe denominator of $\frac{25}{1 - ab}$ ends up being $\frac{1}{6}$, so multiplying both numerator and denominator by 6 gives $\boxed{150}$.

 -ThisUsernameIsTaken

 


 Remark: The quadratic need not be solved. The value of $ab$ can be found through Vieta's.

 




	\item (\textit{AIME I 1986}) Determine $3x_4+2x_5$ if $x_1$, $x_2$, $x_3$, $x_4$, and $x_5$ satisfy the system of equations below.


 $2x_1+x_2+x_3+x_4+x_5=6$$x_1+2x_2+x_3+x_4+x_5=12$$x_1+x_2+2x_3+x_4+x_5=24$$x_1+x_2+x_3+2x_4+x_5=48$$x_1+x_2+x_3+x_4+2x_5=96$
 \textbf{Answer: }181



\textbf{Solution}

Adding all five \href{https://artofproblemsolving.com/wiki/index.php?title=Equation}{equations} gives us $6(x_1 + x_2 + x_3 + x_4 + x_5) = 6(1 + 2 + 4 + 8 + 16)$ so $x_1 + x_2 + x_3 + x_4 + x_5 = 31$.  Subtracting this from the fourth given equation gives $x_4 = 17$ and subtracting it from the fifth given equation gives $x_5 = 65$, so our answer is $3\cdot17 + 2\cdot65 = \boxed{181}$.

 





\textbf{Solution 2}

Subtracting the first equation from every one of the other equations yields
 \begin{align*} x_2-x_1&=6\\ x_3-x_1&=18\\ x_4-x_1&=42\\ x_5-x_1&=90 \end{align*}
Thus
 \begin{align*} 2x_1+x_2+x_3+x_4+x_5&=6\\ 2x_1+(x_1+6)+(x_1+18)+(x_1+42)+(x_1+90)&=6\\ 6x_1+156&=6\\ x_1&=-25 \end{align*}
Using the previous equations,
 \[3x_4+2x_5=3(x_1+42)+2(x_1+90)=\boxed{181}\]

 




	\item (\textit{AIME I 1986}) What is that largest \href{https://artofproblemsolving.com/wiki/index.php?title=Positive\_integer}{positive integer} $n$ for which $n^3+100$ is \href{https://artofproblemsolving.com/wiki/index.php?title=Divisible}{divisible} by $n+10$?


 
 \textbf{Answer: }890



\textbf{Solution 1}

If $n+10 \mid n^3+100$, $\gcd(n^3+100,n+10)=n+10$. Using the \href{https://artofproblemsolving.com/wiki/index.php?title=Euclidean\_algorithm}{Euclidean algorithm}, we have $\gcd(n^3+100,n+10)= \gcd(-10n^2+100,n+10)$ $= \gcd(100n+100,n+10)$ $= \gcd(-900,n+10)$, so $n+10$ must divide $900$. The greatest \href{https://artofproblemsolving.com/wiki/index.php?title=Integer}{integer} $n$ for which $n+10$ divides $900$ is $\boxed{890}$; we can double-check manually and we find that indeed $900\mid 890^3+100$.

 





\textbf{Solution 2 (Simple)}

Let $n+10=k$, then $n=k-10$. Then $n^3+100 = k^3-30k^2+300k-900$ Therefore, $900$ must be divisible by $k$, which is largest when $k=900$ and $n=\boxed{890}$

 





\textbf{Solution 3}

In a similar manner, we can apply synthetic division. We are looking for $\frac{n^3 + 100}{n + 10} = n^2 - 10n + 100 - \frac{900}{n + 10}$. Again, $n + 10$ must be a factor of $900 \Longrightarrow n = \boxed{890}$.

 





\textbf{Solution 4}

The key to this problem is to realize that $n+10 \mid n^3 +1000$ for all $n$. Since we are asked to find the maximum possible $n$ such that $n+10 \mid n^3 +100$, we have: $n+10 \mid ((n^3 +1000) - (n^3 +100) \longrightarrow n+10 \mid 900$. This is because of the property that states that if $a \mid b$ and $a \mid c$, then $a \mid b \pm c$. Since, the largest factor of 900 is itself we have: $n+10=900 \Longrightarrow \boxed{n = 890}$

 





\textbf{Solution 5 (Easy Modular Arithmetic)}

Notice that $n\equiv -10 \pmod{n+10}$. Therefore 
 \[0 \equiv n^3+100\equiv(-10)^3+100=-900 \pmod{n+10} \Rightarrow n+10 | 900 \Rightarrow \max_{n+10 | n^3+100} {n} = \boxed{890}.\]

 





\textbf{Solution 6 (Easiest Factor)}

In the problem, we can see that $n^3 + 100$ has a cube. Immediately think of sum of cubes, which leads us to add 900, resulting in $(n^3 + 10^3) - 900$ --> $(n+10)(n^2 - 10n + 100) - 900$. Since the $(n+10)(n^2 - 10n + 100)$ contains $(n+10)$, this part is divisible by $n+10$ --> therefore the 900 must be divisible by $n+10$. The largest $n$ that satisfies this is $900-10 = \boxed{890}$.

 




	\item (\textit{AIME I 1986}) The pages of a book are numbered $1$ through $n$. When the page numbers of the book were added, one of the page numbers was mistakenly added twice, resulting in an incorrect sum of $1986$. What was the number of the page that was added twice?


 
 \textbf{Answer: }033



\textbf{Solution 1}

Denote the page number as $x$, with $x < n$. The sum formula shows that $\frac{n(n + 1)}{2} + x = 1986$. Since $x$ cannot be very large, disregard it for now and solve $\frac{n(n+1)}{2} = 1986$. The positive root for $n \approx \sqrt{3972} \approx 63$. Quickly testing, we find that $63$ is too large, but if we plug in $62$ we find that our answer is $\frac{62(63)}{2} + x = 1986 \Longrightarrow x = \boxed{033}$.

 





\textbf{Solution 2}

Use the same method as above where you represent the sum of integers from $1$ to $n$ expressed as $\frac{n(n + 1)}{2}$, plus the additional page number $k$. We can establish an upper and lower bound for the number of pages contained in the book, where the upper bound maximizes the sum of the pages and the lower bound maximizes the extra page (and thus minimizes the total sum of the pages.) 

 $\frac{n(n + 1)}{2} = 1986$ is the quadratic you must solve to obtain the upper bound of $n$ and $\frac{n(n + 1)}{2} + n = 1986$ is the quadratic you must solve to obtain the lower bound of $n$. 

 Solving the two equations gives values that are respectively around $62.5$ and $61.5$ with the quadratic formula, and the only integer between the two is $62$. 

 This implies that we can plug in $62$ and come to the same conclusion as the above solution where $x = \boxed{033}$.

 




	\item (\textit{AIME I 1986}) The increasing sequence $1,3,4,9,10,12,13\cdots$ consists of all those positive integers which are \href{https://artofproblemsolving.com/wiki/index.php?title=Powers&action=edit&redlink=1}{powers} of 3 or sums of distinct powers of 3. Find the $100^{\mbox{th}}$ term of this sequence.


 
 \textbf{Answer: }981



\textbf{Solution 1}

Rewrite all of the terms in base 3. Since the numbers are sums of \textit{distinct} powers of 3, in base 3 each number is a sequence of 1s and 0s (if there is a 2, then it is no longer the sum of distinct powers of 3). Therefore, we can recast this into base 2 (binary) in order to determine the 100th number. $100$ is equal to $64 + 32 + 4$, so in binary form we get $1100100$. However, we must change it back to base 10 for the answer, which is $3^6 + 3^5 + 3^2 = 729 + 243 + 9 = \boxed {981}$.

 





\textbf{Solution 2}

Notice that the first term of the sequence is $1$, the second is $3$, the fourth is $9$, and so on. Thus the $64th$ term of the sequence is $729$. Now out of $64$ terms which are of the form $729$ + $'''S'''$, $32$ of them include $243$ and $32$ do not. The smallest term that includes $243$, i.e. $972$, is greater than the largest term which does not, or $854$. So the $96$th term will be $972$, then $973$, then $975$, then $976$, and finally $\boxed{981}$

 





\textbf{Solution 3}

After the $n$th power of 3 in the sequence, the number of terms after that power but before the $(n+1)$th power of 3 is equal to the number of terms before the $n$th power, because those terms after the $n$th power are just the $n$th power plus all the distinct combinations of powers of 3 before it, which is just all the terms before it. Adding the powers of $3$ and the terms that come after them, we see that the $100$th term is after $729$, which is the $64$th term. Also, note that the $k$th term after the $n$th power of 3 is equal to the power plus the $k$th term in the entire sequence. Thus, the $100$th term is $729$ plus the $36$th term. Using the same logic, the $36$th term is $243$ plus the $4$th term, $9$. We now have $729+243+9=\boxed{981}$

 





\textbf{Solution 4}

Writing out a few terms of the sequence until we reach the next power of 3 (27), we see that the $2^{nth}$ term is equal to $3^n$. From here, we can ballpark the range of the 100th term. The 64th term is $3^6$ = $729$ and the 128th term is $3^7$ = $2187$. Writing out more terms of the sequence until the next power of 3 again (81) we can see that the ($2^n$+$2^{n+1}$)/2 term is equal to $3^n$ + $3^{n-1}$. From here, we know that the 96th term is $3^6$ + $3^5$ = $972$. From here, we can construct the 100th term by following the sequence in increasing order. The 97th term is $972 + 1 = 973$, the 98th term is $972 + 3 = 975$, the 99th term is $972 + 3 + 1 = 976$, and finally the 100th term is $972 + 9 = \boxed{981}$

 





\textbf{Solution 5}

The number of terms $3^n$ produces includes each power of 3 ($1, 3^1, ..., 3^n$), the sums of two power of 3s(ex. $3^1 + 1$), three power of 3s (ex. $3^1 + 1 + 3^n$), all the way to the sum of them all. Since there are $n+1$ powers of 3, the one number sum gives us ${n+1\choose 1}$ terms, the two number ${n+1\choose 2}$ terms, all the way to the sum of all the powers which gives us ${n+1\choose n+1}$ terms. Summing all these up gives us $2^{n+1} - 1 ^ {*}$ according to the theorem 

 $\sum_{k=0}^N{{N}\choose{k}} = 2^N$.Since $2^6$ is the greatest power $<100$, then $n=5$ and the sequence would look like {$3^0, ..., 3^5$}, where $3^5$ or $243$ would be the $2^5 - 1 = 63$rd number. The next largest power $729$ would be the 64th number. However, its terms contributed extends beyond 100, so we break it to smaller pieces. Noting that $729$ plus any combination of lower powers ${1, 3^1 . . .3^4}$ is < $729 + 243$, so we can add all those terms($2^5 - 1 = 31$) into our sequence:

 {$3^0, ..., 3^5, 729, 729 + 1, .... 729 + (1 + 3^1 + . . . +3^4)$}Our sequence now has $63 + 1 + 31 = 95$ terms. The remaining $5$ would just be the smallest sums starting with $729 + 243$ or $972$:

 [$972$, $972 + 1$, $972 + 3$, $972+1+3$, $972 + 9$]Hence the 100th term would be $972 + 9 = \boxed{981}$. ~SoilMilk

 ${}^{*}$Note that there isn't a ${n+1\choose 0}$ as choosing 0 numbers will not give you a term.

 


 




	\item (\textit{AIME I 1986}) Let $S$ be the sum of the base $10$ logarithms of all the proper divisors of $1{,}000{,}000$. What is the integer nearest to $S$?


 
 \textbf{Answer: }141



\textbf{Solution 1}

The \href{https://artofproblemsolving.com/wiki/index.php?title=Prime\_factorization}{prime factorization} of $1000000 = 2^65^6$, so there are $(6 + 1)(6 + 1) = 49$ divisors, of which $48$ are proper. The sum of multiple logarithms of the same base is equal to the logarithm of the products of the numbers. 

 Writing out the first few terms, we see that the answer is equal to 
 \[\log 1 + \log 2 + \log 4 + \log 5 +\ldots + \log 1000000 = \log (2^05^0)(2^15^0)(2^25^0)\cdots (2^65^6).\] Each power of $2$ appears $7$ times; and the same goes for $5$. So the overall power of $2$ and $5$ is $7(1+2+3+4+5+6) = 7 \cdot 21 = 147$. However, since the question asks for proper divisors, we exclude $2^65^6$, so each power is actually $141$ times. The answer is thus $S = \log 2^{141}5^{141} = \log 10^{141} = \boxed{141}$.

 \textbf{Simplification} The formula for the product of the divisors of $n$ is $n^{(d(n))/2}$, where $d(n)$ is the number of divisor of $n$.We know that $\log_{10} a + \log_{10} b + \log_{10} c + \log_{10} d...$ and so on equals $\log_{10} (abcd...)$ by sum-product rule of logs, so the problem is reduced to finding the logarithm base 10 of the product of the proper divisors of $10^6$. The product of the divisors, by the earlier formula, is ${(10^6)}^{49/2} = 10^{49*3}$, and since we need the product of only the proper divisors, which means the divisors NOT including the number, $10^6$, itself, we divide $10^{(49*3)}$ by $10^6$ to get $10^{(49*3-6)} = 10^{(141)}$. The base-10 logarithm of this value, in base 10, is clearly $\boxed{141}$

 





\textbf{Solution 2}

Since the prime factorization of $10^6$ is $2^6 \cdot 5^6$, the number of factors in $10^6$ is $7 \cdot 7=49$. You can pair them up into groups of two so each group multiplies to $10^6$. Note that $\log n+\log{(10^6/n)}=\log{n}+\log{10^6}-\log{n}=6$. Thus, the sum of the logs of the divisors is half the number of divisors of $10^6 \cdot 6 -6$ (since they are asking only for proper divisors), and the answer is $(49/2)\cdot 6-6=\boxed{141}$.

 


 


 





\textbf{Solution 3}

Note that we can just pair terms up such that the product is $10^{6}.$ Now, however, note that $10^{3}$ is not included. Therefore we first exclude. We have $\displaystyle\frac{49-1}{2} = 24$ pairs that all multiply to $10^{6}.$ Now we include $10^{3}$ so our current product is $24 \cdot 6 - 3.$ However we dont want to include $10^6$ since we are considering proper factors only so the final answer is $144 + 3 - 6 = \boxed{141}.$

 




	\item (\textit{AIME I 1986}) In $\triangle ABC$, $AB= 425$, $BC=450$, and $AC=510$. An interior point $P$ is then drawn, and segments are drawn through $P$ parallel to the sides of the triangle. If these three segments are of an equal length $d$, find $d$.


 
 \textbf{Answer: }306



\textbf{Solution 1}


                    \begin{figure}[H]
                        \centering
                        \includegraphics[width=0.5\textwidth]{Solution/AIME_I_1986_9_Solution_Figure_1.png}
                    \end{figure}
                    Let the points at which the segments hit the triangle be called $D, D', E, E', F, F'$ as shown above. As a result of the lines being parallel, all three smaller triangles and the larger triangle are \href{https://artofproblemsolving.com/wiki/index.php?title=Similar\_triangles}{similar} ($\triangle ABC \sim \triangle DPD' \sim \triangle PEE' \sim \triangle F'PF$). The remaining three sections are \href{https://artofproblemsolving.com/wiki/index.php?title=Parallelogram}{parallelograms}.

 By similar triangles, $BE'=\frac{d}{510}\cdot450=\frac{15}{17}d$ and $EC=\frac{d}{425}\cdot450=\frac{18}{17}d$. Since $FD'=BC-EE'$, we have $900-\frac{33}{17}d=d$, so $d=\boxed{306}$.

 





\textbf{Solution 2}


                    \begin{figure}[H]
                        \centering
                        \includegraphics[width=0.5\textwidth]{Solution/AIME_I_1986_9_Solution_Figure_2.png}
                    \end{figure}
                    

 Construct cevians $AX$, $BY$ and $CZ$ through $P$.  Place masses of $x,y,z$ on $A$, $B$ and $C$ respectively; then $P$ has mass $x+y+z$.

 Notice that $Z$ has mass $x+y$.  On the other hand, by similar triangles, $\frac{CP}{CZ} = \frac{d}{AB}$.  Hence by mass points we find that 
 \[\frac{x+y}{x+y+z} = \frac{d}{AB}\] Similarly, we obtain 
 \[\frac{y+z}{x+y+z} = \frac{d}{BC} \qquad \text{and} \qquad \frac{z+x}{x+y+z} = \frac{d}{CA}\] Summing these three equations yields 
 \[\frac{d}{AB} + \frac{d}{BC} + \frac{d}{CA} =  \frac{x+y}{x+y+z} +  \frac{y+z}{x+y+z} +  \frac{z+x}{x+y+z} = \frac{2x+2y+2z}{x+y+z} = 2\]

 Hence, 

 $d = \frac{2}{\frac{1}{AB} + \frac{1}{BC} + \frac{1}{CA}} = \frac{2}{\frac{1}{510} + \frac{1}{450} + \frac{1}{425}} = \frac{10}{\frac{1}{85}+\frac{1}{90}+\frac{1}{102}}$$= \frac{10}{\frac{1}{5}\left(\frac{1}{17}+\frac{1}{18}\right)+\frac{1}{102}}=\frac{10}{\frac{1}{5}\cdot\frac{35}{306}+\frac{3}{306}}=\frac{10}{\frac{10}{306}} = \boxed{306}$





\textbf{Solution 3}


                    \begin{figure}[H]
                        \centering
                        \includegraphics[width=0.5\textwidth]{Solution/AIME_I_1986_9_Solution_Figure_3.png}
                    \end{figure}
                    Let the points at which the segments hit the triangle be called $D, D', E, E', F, F'$ as shown above. As a result of the lines being parallel, all three smaller triangles and the larger triangle are \href{https://artofproblemsolving.com/wiki/index.php?title=Similar\_triangles}{similar} ($\triangle ABC \sim \triangle DD'P \sim \triangle PEE' \sim \triangle F'PF$). The remaining three sections are \href{https://artofproblemsolving.com/wiki/index.php?title=Parallelogram}{parallelograms}.

 Since $PDAF'$ is a parallelogram, we find $PF' = AD$, and similarly $PE = BD'$. So $d = PF' + PE = AD + BD' = 425 - DD'$. Thus $DD' = 425 - d$. By the same logic, $EE' = 450 - d$. 

 Since $\triangle DPD' \sim \triangle ABC$, we have the \href{https://artofproblemsolving.com/wiki/index.php?title=Proportion}{proportion}:

 $\frac{425-d}{425} = \frac{PD}{510} \Longrightarrow PD = 510 - \frac{510}{425}d = 510 - \frac{6}{5}d$Doing the same with $\triangle PEE'$, we find that $PE' =510 - \frac{17}{15}d$. Now, $d = PD + PE' = 510 - \frac{6}{5}d + 510 - \frac{17}{15}d \Longrightarrow d\left(\frac{50}{15}\right) = 1020 \Longrightarrow d = \boxed{306}$.

 





\textbf{Solution 4}

Define the points the same as above.

 Let $[CE'PF] = a$, $[E'EP] = b$, $[BEPD'] = c$, $[D'PD] = d$, $[DAF'P] = e$ and $[F'D'P] = f$

 The key theorem we apply here is that the ratio of the areas of 2 similar triangles is the ratio of a pair of corresponding sides squared.

 Let the length of the segment be $x$ and the area of the triangle be $A$, using the theorem, we get:

 $\frac {d + e + f}{A} = \left(\frac {x}{BC}\right)^2$, $\frac {b + c + d}{A}= \left(\frac {x}{AC}\right)^2$, $\frac {a + b + f}{A} = \left(\frac {x}{AB}\right)^2$.Adding all these together and using $a + b + c + d + e + f = A$ we get$\frac {f + d + b}{A} + 1 = x^2 \cdot \left(\frac {1}{BC^2} + \frac {1}{AC^2} + \frac {1}{AB^2}\right)$

 Using \href{https://artofproblemsolving.com/wiki/index.php?title=Corresponding\_angles&action=edit&redlink=1}{corresponding angles} from parallel lines, it is easy to show that $\triangle ABC \sim \triangle F'PF$; since $ADPF'$ and $CFPE'$ are parallelograms, it is easy to show that $FF' = AC - x$

 Now we have the side length \href{https://artofproblemsolving.com/wiki/index.php?title=Ratio}{ratio}, so we have the area ratio$\frac {f}{A} = \left(\frac {AC - x}{AC}\right)^2 = \left(1 - \frac {x}{AC}\right)^2$. By symmetry, we have$\frac {d}{A} = \left(1 - \frac {x}{AB}\right)^2$ and $\frac {b}{A} = \left(1 - \frac {x}{BC}\right)^2$

 Substituting these into our initial equation, we have$1 + \sum_{cyc}\left(1 - \frac {x}{AB}\right) - \frac {x^2}{AB^2} = 0$$\implies 1 + \sum_{cyc}1 - 2 \cdot \frac {x}{AB} = 0$$\implies \frac {2}{\frac {1}{AB} + \frac {1}{BC} + \frac {1}{CA}} = x$and the answer follows after some hideous computation.

 





\textbf{Solution 5}

Refer to the diagram in solution 2; let $a^2=[E'EP]$, $b^2=[D'DP]$, and $c^2=[F'FP]$. Now, note that $[E'BD]$, $[D'DP]$, and $[E'EP]$ are similar, so through some similarities we find that $\frac{E'P}{PD}=\frac{a}{b}\implies\frac{E'D}{PD}=\frac{a+b}{b}\implies[E'BD]=b^2\left(\frac{a+b}{b}\right)^2=(a+b)^2$. Similarly, we find that $[D'AF]=(b+c)^2$ and $[F'CE]=(c+a)^2$, so $[ABC]=(a+b+c)^2$. Now, again from similarity, it follows that $\frac{d}{510}=\frac{a+b}{a+b+c}$, $\frac{d}{450}=\frac{b+c}{a+b+c}$, and $\frac{d}{425}=\frac{c+a}{a+b+c}$, so adding these together, simplifying, and solving gives $d=\frac{2}{\frac{1}{425}+\frac{1}{450}+\frac{1}{510}}=\frac{10}{\frac{1}{85}+\frac{1}{90}+\frac{1}{102}}=\frac{10}{\frac{1}{5}\left(\frac{1}{17}+\frac{1}{18}\right)+\frac{1}{102}}=\frac{10}{\frac{1}{5}\cdot\frac{35}{306}+\frac{3}{306}}$$=\frac{10}{\frac{10}{306}}=\boxed{306}$.

 





\textbf{Solution 6}


                    \begin{figure}[H]
                        \centering
                        \includegraphics[width=0.5\textwidth]{Solution/AIME_I_1986_9_Solution_Figure_1.png}
                    \end{figure}
                    Refer to the diagram above. Notice that because $CE'PF$, $AF'PD$, and $BD'PE$ are parallelograms, $\overline{DD'} = 425-d$, $\overline{EE'} = 450-d$, and $\overline{FF'} = 510-d$. 

 Let $F'P = x$. Then, because $\triangle ABC \sim \triangle F'PF$, $\frac{AB}{AC}=\frac{F'P}{F'F}$, so $\frac{425}{510}=\frac{x}{510-d}$. Simplifying the LHS and cross-multiplying, we have $6x=2550-5d$. From the same triangles, we can find that $FP=\frac{18}{17}x$.

 $\triangle PEE'$ is also similar to $\triangle F'PF$. Since $EF'=d$, $EP=d-x$. We now have $\frac{PE}{EE'}=\frac{F'P}{FP}$, and $\frac{d-x}{450-d}=\frac{17}{18}$. Cross multiplying, we have $18d-18x=450 \cdot 17-17d$. Using the previous equation to substitute for $x$, we have: 
 \[18d-3\cdot2550+15d=450\cdot17-17d\] This is a linear equation in one variable, and we can solve to get $d=\boxed{306}$

 
                \begin{itemize}
                    \item I did not show the multiplication in the last equation because most of it cancels out when solving.
                \end{itemize}


                (Note: I chose $F'P$ to be $x$ only because that is what I had written when originally solving. The solution would work with other choices for $x$.)

 




	\item (\textit{AIME I 1986}) In a parlor game, the magician asks one of the participants to think of a three digit number $(abc)$ where $a$, $b$, and $c$ represent digits in base $10$ in the order indicated. The magician then asks this person to form the numbers $(acb)$, $(bca)$, $(bac)$, $(cab)$, and $(cba)$, to add these five numbers, and to reveal their sum, $N$. If told the value of $N$, the magician can identify the original number, $(abc)$. Play the role of the magician and determine the $(abc)$ if $N= 3194$.


 
 \textbf{Answer: }358



\textbf{Solution 1}

Let $m$ be the number $100a+10b+c$. Observe that $3194+m=222(a+b+c)$ so

 
 \[m\equiv -3194\equiv -86\equiv 136\pmod{222}\]

 This reduces $m$ to one of $136, 358, 580, 802$. But also $a+b+c=\frac{3194+m}{222}>\frac{3194}{222}>14$ so $a+b+c\geq 15$. Recall that $a, b, c$ refer to the digits the three digit number $(abc)$, so of the four options, only $m = \boxed{358}$ satisfies this inequality.

 





\textbf{Solution 2}

As in Solution 1, $3194 + m \equiv 222(a+b+c) \pmod{222}$, and so as above we get $m \equiv 136 \pmod{222}$.We can also take this equation modulo $9$; note that $m \equiv a+b+c \pmod{9}$, so

 
 \[3194 + m \equiv 222m \implies 5m \equiv 8 \implies m \equiv 7 \pmod{9}.\]

 Therefore $m$ is $7$ mod $9$ and $136$ mod $222$. There is a shared factor in $3$ in both, but the Chinese Remainder Theorem still tells us the value of $m$ mod $666$, namely $m \equiv 358$ mod $666$. We see that there are no other 3-digit integers that are $358$ mod $666$, so $m = \boxed{358}$.

 


 





\textbf{Solution 3}

Let $n=abc$ then
 \[N=222(a+b+c)-n\]
 \[N=222(a+b+c)-100a-10b-c=3194\]Since $0<100a+10b+c<1000$, we get the inequality
 \[N<222(a+b+c)<N+1000\]
 \[3194<222(a+b+c)<4194\]
 \[14<a+b+c<19\]Checking each of the multiples of $222$ from $15\cdot222$ to $18\cdot222$ by subtracting $N$ from each $222(a+b+c)$, we quickly find $n=\boxed{358}$

 





\textbf{Solution 4}

The sum of the five numbers is $222(a+b+c)-100a-10b-c=122a+212b+221c=122(a+b+c)+9(10b+11c)=3194$ We can see that $3194 \equiv 8$ (mod $9$) and $122 \equiv 5$ (mod $9$) so we need to make sure that $a+b+c \equiv 7$ (mod $9$) by some testing. So we let $a+b+c=9k+7$

 Then, we know that $1\leq a+b+c \leq 27$ so only $7,16,25$ lie in the interval

 When we test $a+b+c=25, 10b+11c=16$, impossible

 When we test $a+b+c=16, 10b+11c=138, b=5,c=8,a=3$

 When we test $a+b+c=7, 10b+11c=260$, well, it's impossible

 The answer is $\boxed{358}$ then

 




	\item (\textit{AIME I 1986}) The polynomial $1-x+x^2-x^3+\cdots+x^{16}-x^{17}$ may be written in the form $a_0+a_1y+a_2y^2+\cdots +a_{16}y^{16}+a_{17}y^{17}$, where $y=x+1$ and the $a_i$'s are constants. Find the value of $a_2$.


 
 \textbf{Answer: }816



\textbf{Solution 1}

Using the \href{https://artofproblemsolving.com/wiki/index.php?title=Geometric\_series}{geometric series} formula, $1 - x + x^2 + \cdots - x^{17} = \frac {1 - x^{18}}{1 + x} = \frac {1-x^{18}}{y}$. Since $x = y - 1$, this becomes $\frac {1-(y - 1)^{18}}{y}$. We want $a_2$, which is the coefficient of the $y^3$ term in $-(y - 1)^{18}$ (because the $y$ in the denominator reduces the degrees in the numerator by $1$). By the \href{https://artofproblemsolving.com/wiki/index.php?title=Binomial\_Theorem}{Binomial Theorem}, this is $(-1) \cdot (-1)^{15}{18 \choose 3} = \boxed{816}$.

 





\textbf{Solution 2}

Again, notice $x = y - 1$. So

 
 \begin{align*}1 - x + x^2 + \cdots - x^{17} & = 1 - (y - 1) + (y - 1)^2 - (y - 1)^3 + \cdots - (y - 1)^{17} \\ & = 1 + (1 - y) + (1 - y)^2 + (1 - y)^3 \cdots + (1 - y)^{17}\end{align*}
.

 We want the coefficient of the $y^2$ term of each power of each binomial, which by the binomial theorem is ${2\choose 2} + {3\choose 2} + \cdots + {17\choose 2}$. The \href{https://artofproblemsolving.com/wiki/index.php?title=Hockey\_Stick\_Identity}{Hockey Stick Identity} tells us that this quantity is equal to ${18\choose 3} = \boxed{816}$.

 





\textbf{Solution 3}

Again, notice $x=y-1$. Substituting $y-1$ for $x$ in $f(x)$ gives: 
 \begin{align*}1 - x + x^2 + \cdots - x^{17} & = 1 - (y - 1) + (y - 1)^2 - (y - 1)^3 + \cdots - (y - 1)^{17} \\ & = 1 + (1 - y) + (1 - y)^2 + (1 - y)^3 \cdots + (1 - y)^{17}\end{align*}
.From binomial theorem, the coefficient of the $y^2$ term is ${2\choose 0} + {3\choose 1} + \cdots + {17\choose 15}$. This is actually the sum of the first 16 triangular numbers, which evaluates to $\frac{(16)(17)(18)}{6} = \boxed{816}$.

 





\textbf{Solution 4(calculus)}

Let $f(x)=1-x+x^2-x^3+\cdots+x^{16}-x^{17}$ and $g(y)=a_0+a_1y+a_2y^2+\cdots +a_{16}y^{16}+a_{17}y^{17}$.

 Then, since $f(x)=g(y)$,
 \[\frac{d^2f}{dx^2}=\frac{d^2g}{dy^2}\]$\frac{d^2f}{dx^2} = 2\cdot 1 - 3\cdot 2x+\cdots-17\cdot 16x^{15}$ by the power rule.

 Similarly, $\frac{d^2g}{dy^2} = a_2(2\cdot 1) + a_3(3\cdot 2y)+\cdots+a_{17}(17\cdot 16y^{15})$

 Now, notice that if $x = -1$, then $y = 0$, so $f^{''}(-1) = g^{''}(0)$

 $g^{''}(0)= 2a_2$, and $f^{''}(-1) = 2\cdot 1 + 3\cdot 2 +\cdots + 16\cdot 17$.

 Now, we can use the hockey stick theorem to see that $2\cdot 1 + 3\cdot 2 +\cdots + 16\cdot 17 = 2\binom{18}{3}$

 Thus, $2a_2 = 2\binom{18}{3}\rightarrow a_2 = \binom{18}{3}=\boxed{816}$

 -AOPS81619

 





\textbf{Solution 5 (Linear Algebra)}

Let $V$ be the vector space of polynomials of degree $\leq 17,$ and let $B = \{1, x, x^2, ..., x^{17} \}$ and $C = \{1, (x+1), (x+1)^2, ..., (x+1)^{17} \}$ be two bases for $V$.Let $\vec{v} \in V$ be the polynomial given in the problem, and it is easy to see that $[ \vec{v} ]_B = \langle 1, -1, 1, -1, ... , 1, -1 \rangle.$

 Note that the transformation matrix from $C$ to $B$ can be easily found to be $P_{C \to B} = [ [\vec{c_1}]_B [\vec{c_2}]_B ... [\vec{c_3}]_B ] = \begin{bmatrix} \tbinom{0}{0} & \tbinom{1}{0} & \tbinom{2}{0} & \cdots & \tbinom{17}{0} \\ 0 & \tbinom{1}{1} & \tbinom{2}{1} & \cdots & \tbinom{17}{1} \\ 0 & 0 & \tbinom{2}{2} & \cdots & \tbinom{17}{2} \\ \vdots & \vdots & \vdots & \ddots & \vdots \\ 0 & 0 & 0 & \cdots & \tbinom{17}{17} \end{bmatrix} .$

 I claim that $P_{B \to C} = \begin{bmatrix} \tbinom{0}{0} & -\tbinom{1}{0} & \tbinom{2}{0} & \cdots & -\tbinom{17}{0} \\ 0 & \tbinom{1}{1} & -\tbinom{2}{1} & \cdots & \tbinom{17}{1} \\ 0 & 0 & \tbinom{2}{2} & \cdots & -\tbinom{17}{2} \\ \vdots & \vdots & \vdots & \ddots & \vdots \\ 0 & 0 & 0 & \cdots & \tbinom{17}{17} \end{bmatrix} ,$ where the term $\dbinom{n}{k}$ is negated if $n+k$ is odd. 

 One can prove that the $i$th row of $P_{C \to B}$ dotted with the $j$th column of $P_{B \to C}$ is $\delta_{i, j}$ by using combinatorial identities, which is left as an exercise for the reader. Thus, since the two matrices multiply to form $\mathbb{I}_{18},$ we have proved that $P_{B \to C} = \begin{bmatrix} \tbinom{0}{0} & -\tbinom{1}{0} & \tbinom{2}{0} & \cdots & -\tbinom{17}{0} \\ 0 & \tbinom{1}{1} & -\tbinom{2}{1} & \cdots & \tbinom{17}{1} \\ 0 & 0 & \tbinom{2}{2} & \cdots & -\tbinom{17}{2} \\ \vdots & \vdots & \vdots & \ddots & \vdots \\ 0 & 0 & 0 & \cdots & \tbinom{17}{17} \end{bmatrix} .$

 To find the coordinates of $\vec{v}$ under basis $C,$ we compute the product $[ \vec{v} ]_C = P_{B \to C} [\vec{v} ]_B = \begin{bmatrix} \tbinom{0}{0} & -\tbinom{1}{0} & \tbinom{2}{0} & \cdots & -\tbinom{17}{0} \\ 0 & \tbinom{1}{1} & -\tbinom{2}{1} & \cdots & \tbinom{17}{1} \\ 0 & 0 & \tbinom{2}{2} & \cdots & -\tbinom{17}{2} \\ \vdots & \vdots & \vdots & \ddots & \vdots \\ 0 & 0 & 0 & \cdots & \tbinom{17}{17} \end{bmatrix} \begin{bmatrix} 1 \\ -1 \\ 1 \\ \vdots \\ -1 \end{bmatrix} = \begin{bmatrix} \sum_{n=0}^{17} \tbinom{n}{0} \\ -\sum_{n=1}^{17} \tbinom{n}{1} \\  \sum_{n=2}^{17} \tbinom{n}{2} \\ \vdots \\  -\sum_{n=17}^{17} \tbinom{n}{17} \end{bmatrix} = \begin{bmatrix} \tbinom{18}{1} \\ - \tbinom{18}{2} \\ \tbinom{18}{3} \\ \vdots \\ -\tbinom{18}{18} \end{bmatrix},$ where the last equality was obtained via Hockey Stick Identity.

 Thus, our answer is $a_2 =  [ [ \vec{v} ]_C ]_3 = \dbinom{18}{3} = \boxed{816}.$

 -fidgetboss\_4000

 




	\item (\textit{AIME I 1986}) Let the sum of a set of numbers be the sum of its elements. Let $S$ be a set of positive integers, none greater than 15. Suppose no two disjoint subsets of $S$ have the same sum. What is the largest sum a set $S$ with these properties can have?


 
 \textbf{Answer: }061



\textbf{Solution 1}

The sum of set \( S \) is \( 15(16)/2 = 120 \).

 Let \( A \) and \( B \) be subsets of \( S \).

 Assume \( A = B \). Then the greatest possible value of \( A \) or \( B \) is \( 120/2 = 60 \).

 We want \( A \neq B : \text{Greedy on}~ A ~\text{does not influence}~ B\), so either \( A = 61 \) and \( B = 59 \) or \( A = 59 \) and \( B=61 \). We want the largest sum, so our answer is $\boxed{061}$.

 





\textbf{Solution 2}

By using the greedy algorithm, we obtain $\boxed{061}$, with $S=\{ 15,14,13,11,8\}$. We must now prove that no such set has sum greater than 61. Suppose such a set $S$ existed. Then $S$ must have more than 4 elements, otherwise its sum would be at most $15+14+13+12=54$.  

 $S$ can't have more than 5 elements. To see why, note that at least $\dbinom{6}{0} + \dbinom{6}{1} + \dbinom{6}{2} + \dbinom{6}{3} + \dbinom{6}{4}=57$ of its subsets have at most four elements (the number of subsets with no elements plus the number of subsets with one element and so on), and each of them have sum at most 54. By the Pigeonhole Principle, two of these subsets would have the same sum. If those subsets were disjoint, we would directly arrive at a contradiction; if not, we could remove the common elements to get two disjoint subsets.

 Thus, $S$ would have to have 5 elements. $S$ contains both 15 and 14 (otherwise its sum is at most $10+11+12+13+15=61$). It follows that $S$ cannot contain both $a$ and $a-1$ for any $a\leq 13$, or the subsets $\{a,14\}$ and $\{a-1,15\}$ would have the same sum. So now $S$ must contain 13 (otherwise its sum is at most $15+14+12+10+8=59$), and $S$ cannot contain 12, or the subsets $\{12,15\}$ and $\{13,14\}$ would have the same sum.

 Now the only way $S$ could have sum at least $62=15+14+13+11+9$ would be if $S=\{ 15,14,13,11,9\}$. But the subsets $\{9,15\}$ and $\{11,13\}$ have the same sum, so this set does not work. Therefore no $S$ with sum greater than 61 is possible and 61 is indeed the maximum.

 


 \href{https://mathproblemsolvingskills.wordpress.com/}{MathProblemSolvingSkills.com}

 




	\item (\textit{AIME I 1986}) In a sequence of coin tosses, one can keep a record of instances in which a tail is immediately followed by a head, a head is immediately followed by a head, and etc. We denote these by TH, HH, and etc. For example, in the sequence TTTHHTHTTTHHTTH of 15 coin tosses we observe that there are two HH, three HT, four TH, and five TT subsequences. How many different sequences of 15 coin tosses will contain exactly two HH, three HT, four TH, and five TT subsequences?


 
 \textbf{Answer: }560



\textbf{Solution}

Let's consider each of the sequences of two coin tosses as an \href{https://artofproblemsolving.com/wiki/index.php?title=Operation}{operation} instead; this operation takes a string and adds the next coin toss on (eg, <tt>THHTH</tt> + <tt>HT</tt> = <tt>THHTHT</tt>). We examine what happens to the last coin toss. Adding <tt>HH</tt> or <tt>TT</tt> is simply an \href{https://artofproblemsolving.com/wiki/index.php?title=Identity}{identity} for the last coin toss, so we will ignore them for now. However, adding <tt>HT</tt> or <tt>TH</tt> switches the last coin. <tt>H</tt> switches to <tt>T</tt> three times, but <tt>T</tt> switches to <tt>H</tt> four times; hence it follows that our string will have a structure of <tt>THTHTHTH</tt>. 

 Now we have to count all of the different ways we can add the identities back in. There are 5 <tt>TT</tt> subsequences, which means that we have to add 5 <tt>T</tt> into the strings, as long as the new <tt>T</tt>s are adjacent to existing <tt>T</tt>s. There are already 4 <tt>T</tt>s in the sequence, and since order doesn’t matter between different tail flips this just becomes the ball-and-urn argument. We want to add 5 balls into 4 urns, which is the same as 3 dividers; hence this gives ${{5+3}\choose3} = 56$ combinations. We do the same with 2 <tt>H</tt>s to get ${{2+3}\choose3} = 10$ combinations; thus there are $56 \cdot 10 = \boxed {560}$ possible sequences.

 \textbf{Slight Variation} The structure of the final order is <tt>T\_H\_T\_H\_T\_H\_T\_H\_</tt>, and there are 4 spots to put the 2 heads in, and 4 spots to put the 5 tails in.  By using the formula for distributing r identical objects into n distinct boxes $\dbinom{r+n-1}{r}$ and multiplication, the answer is  ${{2+4-1}\choose2} \cdot{{5+4-1}\choose5}=560$ ~Slight edits in LaTeX by EthanSpoon

 


 





\textbf{Solution 2}

We need exactly four TH subsequences, so we can place additional T's and H's between four blocks of TH:

 \_TH\_TH\_TH\_TH\_

 Since we cannot put additional TH sequences between the four blocks, each of the three spaces in between them will contain a HT subsequence. So now we just need two HH sequences and five TT sequences. 

 For these criteria to be satisfied, we need at least $2$ more H's and $5$ more T's. Since we have a total of $7$ additional flips left, we will have exactly $2$ more H's and $5$ more T's. For reasons explained below, the T's can be  distributed to the first four spaces and the H's can be distributed to the last 4* (since we cannot have anymore TH sequences, the T's are always to the right of the H's). By stars and bars, there are 
 \[{8 \choose 3}{5 \choose 2}=(56)(10)=560\]ways do distribute the T's and H's. Hence, there are $\boxed {560}$ different sequences. 

 
                \begin{itemize}
                    \item [We want to divide $5$ T's among the spaces to form exactly five TT subsequences. Each time we divide the group of T's (TTTTT) to make an additional group, we lose one TT sequence. However, each of the first 4 spaces are adjacent to one T. So each time we divide the T's, if we place the group in one of the first four spaces, we gain back another TT sequence. We start with one group that has four TT sequences, so when we place them in the spaces, we get $1+4=5$ TT sequences. Since the number of TT sequences stays the same each time we divide the T's into smaller groups, we will always have $5$ TT sequences if we divide the additional T's among the first $4$ spaces. We can use identical logic to show why we can divide the two H's among the last four spaces to get exactly two HH subsequences.]
                \end{itemize}


                




	\item (\textit{AIME I 1986}) The shortest distances between an interior \href{https://artofproblemsolving.com/wiki/index.php?title=Diagonal}{diagonal} of a rectangular \href{https://artofproblemsolving.com/wiki/index.php?title=Parallelepiped}{parallelepiped}, $P$, and the edges it does not meet are $2\sqrt{5}$, $\frac{30}{\sqrt{13}}$, and $\frac{15}{\sqrt{10}}$. Determine the \href{https://artofproblemsolving.com/wiki/index.php?title=Volume}{volume} of $P$.


 
 \textbf{Answer: }750



\textbf{Solution}


                    \begin{figure}[H]
                        \centering
                        \includegraphics[width=0.5\textwidth]{Solution/AIME_I_1986_14_Solution_Figure_1.png}
                    \end{figure}
                    In the above diagram, we focus on the line that appears closest and is parallel to $BC$. All the blue lines are perpendicular lines to $BC$ and their other points are on $AB$, the main diagonal. The green lines are projections of the blue lines onto the bottom face; all of the green lines originate in the corner and reach out to $AC$, and have the same lengths as their corresponding blue lines. So we want to find the shortest distance between $AC$ and that corner, which is $\frac {wl}{\sqrt {w^2 + l^2}}$.

 So we have:
 \[\frac {lw}{\sqrt {l^2 + w^2}} = \frac {10}{\sqrt {5}}\]
 \[\frac {hw}{\sqrt {h^2 + w^2}} = \frac {30}{\sqrt {13}}\]
 \[\frac {hl}{\sqrt {h^2 + l^2}} = \frac {15}{\sqrt {10}}\]

 Notice the familiar roots: $\sqrt {5}$, $\sqrt {13}$, $\sqrt {10}$, which are $\sqrt {1^2 + 2^2}$, $\sqrt {2^2 + 3^2}$, $\sqrt {1^2 + 3^2}$, respectively. (This would give us the guess that the sides are of the ratio 1:2:3, but let's provide the complete solution.)

 
 \[\frac {l^2w^2}{l^2 + w^2} = \frac {1}{\frac {1}{l^2} + \frac {1}{w^2}} = 20\]
 \[\frac {h^2w^2}{h^2 + w^2} = \frac {1}{\frac {1}{h^2} + \frac {1}{w^2}} = \frac {900}{13}\]
 \[\frac {h^2l^2}{h^2 + l^2} = \frac {1}{\frac {1}{h^2} + \frac {1}{l^2}} = \frac {45}{2}\]

 We invert the above equations to get a system of linear equations in $\frac {1}{h^2}$, $\frac {1}{l^2}$, and $\frac {1}{w^2}$:
 \[\frac {1}{l^2} + \frac {1}{w^2} = \frac {45}{900}\]
 \[\frac {1}{h^2} + \frac {1}{w^2} = \frac {13}{900}\]
 \[\frac {1}{h^2} + \frac {1}{l^2} = \frac {40}{900}\]

 We see that $h = 15$, $l = 5$, $w = 10$. Therefore $V = 5 \cdot 10 \cdot 15 = \boxed{750}$

 




	\item (\textit{AIME I 1986}) Let triangle $ABC$ be a right triangle in the $xy$-plane with a right angle at $C_{}$. Given that the length of the hypotenuse $AB$ is $60$, and that the medians through $A$ and $B$ lie along the lines $y=x+3$ and $y=2x+4$ respectively, find the area of triangle $ABC$.


 
 \textbf{Answer: }400



\textbf{Solution 1}

Let $\theta_1$ be the angle that the median through $A$ makes with the positive $y$-axis, and let $\theta_2$ be the angle that the median through $B$ makes with the positive $x$-axis.  The tangents of these two angles are the slopes of the respective medians; in other words, $\tan \theta_1 = 1$, and $\tan \theta_2 =2$.

 Let $\theta$ be the angle between the medians through $A$ and $B$.  Then by the \href{https://artofproblemsolving.com/wiki/index.php?title=Trigonometric\_identities\#Angle\_addition\_identities}{ tangent angle subtraction formula},
 \[\tan \theta = \tan (\theta_2 - \theta_1) = \frac{\tan \theta_2 - \tan \theta_1}{1 + \tan \theta_1 \tan \theta_2} = \frac{2-1}{1 + 2 \cdot 1 } = \frac{1}{3}.\]

 
                    \begin{figure}[H]
                        \centering
                        \includegraphics[width=0.5\textwidth]{Solution/AIME_I_1986_15_Solution_Figure_1.png}
                    \end{figure}
                    Let $M$ be the midpoint of $BC,$ let $N$ be the midpoint of $AC,$ and let $G$ be the intersection of these two medians.  Let $BC = a,$ and let $AC = b$; let $\alpha = \angle MAC$, and let $\beta = \angle BNC$.

 With this setup, $\theta = \beta - \alpha$.  The reason is that $\beta$ is an exterior angle of $\triangle AGN$, and thus $\beta = \alpha + \angle AGN = \alpha + \theta$.

 We also know that $\tan \alpha = \dfrac{(1/2) a}{b} = \dfrac{a}{2b}$, and $\tan \beta = \dfrac{a}{(1/2)b} = \dfrac{2a}{b}$.  Then
 \begin{align*} \frac{1}{3} &= \tan \theta \\ &= \tan (\beta - \alpha) \\ &= \frac{\tan \beta - \tan \alpha}{1 + \tan \beta \tan \alpha} \\ &= \frac{2 \frac{a}{b} - \frac{1}{2} \cdot \frac{a}{b}}{1 + (a/b)^2} \\ &= \frac{(3/2) \cdot (a/b)}{1 + (a/b)^2} . \end{align*}
Multiplying numerator and denominator by $b^2,$ we learn that
 \[\frac{1}{3} = \frac{(3/2) ab}{a^2 + b^2} = \frac{(3/2) ab}{60^2},\]since the hypotenuse of $\triangle ABC$ has length 60.  Solving for $\frac{1}{2}ab$, we find that
 \[\frac{1}{2} ab = \frac{1}{9} \cdot 60^2 = (60/3)^2 = 20^2 = 400.\]~minor edits by EaZ\_Shadow

 





\textbf{Solution 2}

Translate so that the \href{https://artofproblemsolving.com/wiki/index.php?title=Median}{medians} are $y = x$, and $y = 2x$, then model the \href{https://artofproblemsolving.com/wiki/index.php?title=Point}{points} $A: (a,a)$ and $B: (b,2b)$. $(0,0)$ is the \href{https://artofproblemsolving.com/wiki/index.php?title=Centroid}{centroid}, and is the average of the vertices, so $C: (- a - b, - a - 2b)$
$AB = 60$ so 

 $3600 = (a - b)^2 + (2b - a)^2$
$3600 = 2a^2 + 5b^2 - 6ab \ \ \ \ (1)$$AC$ and $BC$ are \href{https://artofproblemsolving.com/wiki/index.php?title=Perpendicular}{perpendicular}, so the product of their \href{https://artofproblemsolving.com/wiki/index.php?title=Slope}{slopes} is $-1$, giving

 $\left(\frac {2a + 2b}{2a + b}\right)\left(\frac {a + 4b}{a + 2b}\right) = - 1$
$2a^2 + 5b^2 = - \frac {15}{2}ab \ \ \ \  (2)$Combining $(1)$ and $(2)$, we get $ab = - \frac {800}{3}$

 Using the \href{https://artofproblemsolving.com/wiki/index.php?title=Determinant}{determinant} product for area of a triangle (this simplifies nicely, add columns 1 and 2, add rows 2 and 3), the area is $\left|\frac {3}{2}ab\right|$, so we get the answer to be $400$.

 
Note: If you don't know the determinant product for the area of a triangle, the shoelace theorem also works. It is basically the same thing but more well-known.

 





\textbf{Solution 3}

The only relevant part about the xy plane here is that the slopes of the medians determine an angle between them that we will use.This solution uses the tangent subtraction identity $\tan(\alpha -\beta)=\frac{\tan\alpha -\tan\beta}{1+\tan\alpha \tan\beta}$. Therefore, the tangent of the acute angle between the medians from A and B will be $\frac{2-1}{1+2 \cdot 1}= \frac{1}{3}$.

 Let $AC =x$ and $BC =y$. Let the midpoint of AC be M and the midpoint of BC be N. Let the centroid be G. By exterior angles, $\angle CMB - \angle CAN = \angle AGM$. However we know that since $\angle AGM$ is the acute angle formed by the medians, $\tan \angle AGM=\frac{1}{3}$. We can express the tangents of the other two angles in terms of $x$ and $y$. $\tan \angle CMB= \frac{y}{\frac{x}{2}}=\frac{2y}{x}$ while $\tan \angle CAN= \frac{\frac{y}{2}}{x}=\frac{y}{2x}$. For simplification, let $\frac{y}{x}=r$.

 By the tangent subtraction identity, $\frac{2r-\frac{r}{2}}{1+2r \cdot \frac{r}{2}}=\frac{3r}{2(1+r^2)}=\frac{1}{3}$.We get the quadratic $2r^2 -9r+2=0$, which solves to give $r=\frac{9 \pm \sqrt{65}}{4}$. It does not matter which one is picked, because the two roots multiply to 1, so switching from one root to another is like switching the lengths of AC and BC.

 We choose the positive root and we can plug $y= \left (\frac{9 + \sqrt{65}}{4} \right )x$ into of course the PYTHAGOREAN THEOREM $x^2 + y^2 =3600$ and solve for $x^2=200(9-\sqrt{65})$. We want the area which is $\frac{xy}{2}=\frac{x^2(\frac{9 + \sqrt{65}}{4})}{2}=\frac{200(9-\sqrt{65})(\frac{9 + \sqrt{65}}{4})}{2}=\frac{200\cdot\frac{16}{4}}{2}=400$.

 





\textbf{Solution 4}

We first seek to find the angle between the lines $y = x + 3$ and $y = 2x + 4$.
                    \begin{figure}[H]
                        \centering
                        \includegraphics[width=0.5\textwidth]{Solution/AIME_I_1986_15_Solution_Figure_2.png}
                    \end{figure}
                    Let the acute angle the red line makes with the $x-$ axis be $\alpha$ and the acute angle the blue line makes with the $x-$ axis be $\beta$. Then, we know that $\tan \alpha = 1$ and $\tan \beta = 2$. Note that the acute angle between the red and blue lines is clearly $\beta - \alpha$. Therefore, we have that:
 \[\tan (\beta - \alpha) = \frac{2 - 1}{1 + 2} = \frac{1}{3}\]It follows that $\cos (\beta - \alpha) = \frac{1}{\sqrt{10}}$ and $\sin (\beta - \alpha) = \frac{3}{\sqrt{10}}$. From now on, refer to $\theta = \beta - \alpha$
                    \begin{figure}[H]
                        \centering
                        \includegraphics[width=0.5\textwidth]{Solution/AIME_I_1986_15_Solution_Figure_3.png}
                    \end{figure}
                    Suppose $ABC$ is our desired triangle. Let $E$ be the midpoint of $CA$ and $D$ be the midpoint of $AB$ such that $CE = EA = m$ and $AD = DB = n$. Let $G$ be the centroid of the triangle (in other words, the intersection of $EB$ and $CD$). It follows that if $CG = 2x$, $GD = x$ and if $GB = 2y$, $GE = y$. By the Law of Cosines on $\triangle GEB$, we get that:
 \[GE^2 + GC^2 - 2GE \cdot GC \cdot \cos \theta = CE^2 \Longleftrightarrow\]
 \[(2x)^2 + y^2 - (2x)(y)(2)\cos \theta = m^2\]By the Law of Cosines on $\triangle GDB$, we get that:
 \[DG^2 + GB^2 -2 DG \cdot GB \cdot \cos \theta = DB^2 \Longleftrightarrow\]
 \[(2y)^2+x^2-(2y)(x)(2)\cos \theta = n^2\]Adding yields that:
 \[5x^2+5y^2 - 8xy\cos \theta = m^2 + n^2\]However, note that $(2m)^2 + (2n)^2 = 60^2 \Longleftrightarrow m^2+n^2 = 30^2$. Therefore,
 \[5x^2+5y^2 - 8xy\cos \theta = 30^2\]We also know that by the Law of Cosines on $\triangle CGB$,
 \[CG^2 + GB^2 - 2CG \cdot GB \cdot \cos (180 - \theta) = CB^2 \Longleftrightarrow\]
 \[(2x)^2 + (2y)^2 + 2(2x)(2y) \cos \theta = 60^2 \Longleftrightarrow\]
 \[x^2+y^2+2xy \cos \theta = 30^2 \Longleftrightarrow\]
 \[5x^2+5y^2 + 10xy \cos \theta = 30^2 \cdot 5\]Subtracting this from the $5x^2+5y^2 - 8xy\cos \theta = 30^2$ we got earlier yields that:
 \[xy \cos \theta = 200\]but recall that $\cos \theta = \frac{3}{\sqrt{10}}$ to get that:
 \[xy = \frac{200 \sqrt{10}}{3}\]Plugging this into the $x^2 + y^2 + 2xy \cos \theta = 30^2$, we get that:
 \[x^2 + y^2 = 500\]Aha! How convenient. Recall that $(2x)^2 + y^2 - (2x)(y)(2)\cos \theta = m^2$ and $(2y)^2+x^2-(2y)(x)(2)\cos \theta = n^2$. Then, we clearly have that:
 \[m^2n^2 = ((2x)^2 + y^2 - (2x)(y)(2)\cos \theta)((2y)^2+x^2-(2y)(x)(2)\cos \theta = n^2) \Longleftrightarrow\]
 \[m^2n^2 = 17x^2y^2 - 4xy \cos \theta (5x^2 + 5y^2) + 16x^2y^2 \cos ^2 \theta + 4x^4 + 4y^4 \Longleftrightarrow\]
 \[m^2n^2 = 17 \cdot \frac{200^2 \cdot 10}{9} - 800(5 \cdot 500) + 16 \cdot 200^2 + 4 \left( (x^2+y^2)^2-2x^2y^2\right) \Longleftrightarrow\]
 \[\frac{m^2n^2}{100^2} = \frac{680}{9} - 8 \cdot 5 \cdot 5 + 16 \cdot 2^2 + 4 \left( 5^2 - 2 \cdot \frac{2^2 \cdot 10}{9}\right) \Longleftrightarrow\]
 \[\frac{m^2n^2}{100^2} = 4\]But note that the area of our triangle is $2m \cdot 2n \cdot \frac {1}{2} = 2mn$. As $mn = 200$, we get a final answer of $\boxed{400}$.

 




\end{enumerate}